\section{作成したシステムの概要}
本節では，作成したシステムの概要について説明する．本システムは，ユーザの能動的な動作と連動した観覧車型中継機デバイスの回転およびレーザの発光を通じて，演奏空間全体に多感覚なフィードバックを与えることで，従来の自動演奏技術における演奏インターフェースでは得られない高い没入感をユーザに提供することを目的として設計されている．
本システムは，指揮デバイス，観覧車型の中継機デバイス，楽器デバイスという３つのデバイスから構成される．
% まず，本システムの全体の構成図を図\ref{zentai}に示す．ユーザが指揮デバイスに取り付けられたボタンを押下すると，同期開始情報が無線通信によって中継機デバイスおよび楽器デバイスへ送信され，中継機デバイスが作動を開始する．
演奏開始後，ユーザがスライド式可変抵抗器を操作したうえで指揮デバイスを振ると，加速度センサが振り動作を検知し，中継機デバイスが

ここからは各デバイスの概要について説明する．

\begin{figure}[h]
    \centering
     \centering
    \includegraphics[width=15cm]{./figs/zentai.png}
    \caption{全体の構成図}
    \label{zentai}
\end{figure}
    
    

\subsection{指揮デバイス}
\subsubsection{概要}
まず，指揮デバイスについて説明する．指揮デバイスのシステム構成図を図\ref{fig:sys_siki}に示す．本デバイスは，演奏全体を統括するデバイスであり，主に以下の３つの機能を有する．
第一の機能は，演奏開始トリガーの送信である．ユーザが指揮デバイスに搭載されたタクトスイッチを押下すると，中継機デバイスと楽器デバイスへ演奏開始トリガーを送信する．さらに，中継機デバイスの回転が開始される．

第二の機能は，楽曲のBPM変更である．ユーザがスライド式可変抵抗器を操作した後，デバイスを振ると新しい設定BPM情報が中継機デバイスへ送信される．これにより，中継機デバイスのモータ回転速度がこの値に追従して変化する．そして，その回転に伴うレーザーの通過光を介して各楽器デバイスに情報が伝達され，BPMの変更される．また，楽曲のBPMは5段階に分けて変更できる．

第三の機能は，楽曲の音量変更である．ユーザがスライド式可変抵抗器を操作した後，デバイスを振ると音量情報が各楽器デバイスへマルチキャスト送信されることで，音量が変更される．また，BPMと同様に音量も5段階に分けて変更できる．ここからは，外部設計と内部設計について説明をする．

\begin{figure}[h] % h:ここ, t:上, b:下, p:独立ページ
    \centering % 中央揃え
    \includegraphics[width=0.8\columnwidth]{figs/siki_sys4.png} % 横幅を列幅の80%に指定
    \caption{指揮デバイスのシステム構成図} % 図の説明
    \label{fig:sys_siki} % 本文中で参照するためのラベル
\end{figure}



\subsubsection{外部設計}
本節では指揮デバイスの外部設計について説明する．まず，本デバイスを構成する部品について表\ref{table:siki}に示す．デジタル3軸加速度センサは指揮デバイスの動きの検知に使用し，スライド式可変抵抗器はBPMおよび音量変更に使用される．
当初はデバイスの電力供給源として$\SI{9}{V}$のアルカリ電池を外部電源として使用することを検討していたが，電池なしでも正常動作が確認できたため，軽量化のためPCとのUSB接続から電力を供給するように変更した．Arduino内蔵の電圧レギュレータにより$\SI{5}{V}$および$\SI{3.3}{V}$ を生成し，各モジュールへ安定した電力を供給する．
また，指揮デバイスの回路図を図\ref{fig:kairo_siki} に示す．加速度センサにはGY-291（ADXL345）を用いており，I2C通信によりArduinoと接続する．CSピンをArduinoの$\SI{3.3}{V}$に接続することでI2Cモードに設定し，さらに SDOピンをGNDに接続することでI2Cアドレスを0x53に固定している．また，SDAピンおよびSCLピンはそれぞれデータ通信線およびクロック信号線としてArduinoのA4，A5ピンに接続している．加えて，加速度が内部設定した閾値を超過した場合には，INT1 ピンから Arduino の D2 ピンへ割り込み信号が出力される．
これにより，指揮デバイスの振り動作を検知し，モータ制御用 Arduino との UDP 通信を開始する．

さらに，演奏開始および終了信号の送信に用いるボタン入力判定回路を実装する．ボタンの誤判定を防止するため10k$\Omega$のプルアップ抵抗を用いた回路を構成する．
ボタンが押されていない状態では，入力ピン D3 はプルアップ抵抗を介して 5 V に接続されるため HIGH 状態となる．
一方，ボタン押下時には回路が GND に短絡され，入力電位が LOW に遷移する．指揮デバイスのプログラムでは D3 ピンの状態変化を常時監視することで，
演奏開始および終了操作を判定する．




\begin{table}[H]
    \centering
    \caption{指揮デバイスの構成部品}
    \label{table:siki}
    \begin{tabular}{lll} \hline
        Arduino UNO R4 WiFi& - & 1 \\ \hline
        スライド式可変抵抗器 10 kΩ &-& 2 \\ \hline
        デジタル3軸加速度センサ & GY-291 & 1 \\ \hline
        タクトスイッチ & DTS-63 & 1 \\ \hline
        抵抗器　10k$\Omega$ & - & 1 \\ \hline
        % 外部電源 9 V &-& 1 \\ \hline
        ジャンパーワイヤー & - & 適宜 \\ \hline
    \end{tabular}
\end{table}



\begin{figure}[htbp] % h:ここ, t:上, b:下, p:独立ページ
        \centering % 中央揃え
        \includegraphics[width=0.8\columnwidth]{figs/kairo_siki.png} % 横幅を列幅の80%に指定
        \caption{指揮デバイスの回路図} % 図の説明
        \label{fig:kairo_siki} % 本文中で参照するためのラベル
\end{figure}
% また，指揮デバイスの電力供給源として$\SI{9}{V}$の乾電池を外部電源として使用する．

%%回路の説明もう少しいる？
%%実際の指揮デバイス
図〜〜〜に実際に作成した指揮デバイスを示す．実際の指揮棒のようにユーザが片手で振って操作できるよう，細長い筒状の外装を採用した．また，内部にはブレッドボードに加速度センサ，スライド式可変抵抗器，タクトスイッチを設置する．デバイスのユーザビリティを向上させるため，ユーザが操作するスライダー部分やボタンは外部へ露出させる構造を採用し，さらにデバイスの先端部に加速度センサを設置することで操作時の重心を安定させる．

%%%%%%%%%ソフトウェア
\subsubsection{内部設計}
本節では，指揮デバイスの内部設計について説明する．
まず，指揮デバイスにおけるシーケンス図を図\ref{fig:jou_siki}に示す．
ユーザが指揮デバイスを上下に振ると加速度センサが動作を検知しArduinoに送信される．この時，値が閾値を超過した場合，
観覧車デバイスと無線通信を開始し観覧車デバイスが送信した演奏情報を基に回転する．
また，ボタンを押下した時，各楽器デバイスへ楽曲の初期BPMおよび音量情報が送信されることで演奏が開始される．

\begin{figure}[htbp] % h:ここ, t:上, b:下, p:独立ページ
        \centering % 中央揃え
        \includegraphics[width=0.9\columnwidth]{figs/siki_jou5.pdf} % 横幅を列幅の90%に指定
        \caption{指揮デバイスのシーケンス図} % 図の説明
        \label{fig:jou_siki} % 本文中で参照するためのラベル
\end{figure}

%%あとは通信？？？へけ？？？







\subsection{中継機デバイス}

\subsection{楽器デバイス}
果たしてデバイスなのか＿＿＿